\section{ALeA: an Adaptive Learning Assistant}
The Adaptive Learning Assistant (ALeA)\cite{ALeA:online} constitutes a semantically driven platform
designed to support personalized education.  Building upon the knowledge graph
infrastructure of MathHub, ALeA integrates three
interrelated modules:
\begin{enumerate}
      \item \textbf{Semantic Knowledge Base}:  Lecture materials are transformed
            into a hierarchically organized graph of concepts through semantic
            annotation.  This representation makes explicit the prerequisite
            relationships among topics, thereby enabling precise identification of
            the knowledge components that a learner must master before progressing
            to more advanced material.
      \item \textbf{Learner Model}:   Interaction data (e.g., quiz responses and practice problems) are continuously harvested from
            students using the system.  The model maps these data onto the semantic
            graph, producing a probabilistic estimate of each learner’s mastery of
            individual concepts and of the overall curriculum. 
            The learner model uses Bloom's Taxonomy, which comprises six dimensions of understanding:
            remembering, understanding, applying, analyzing, evaluating, and creating.
            Every property is represented as a probability value.
            The learner model constitutes the most significant component of ALeA for the purposes of this thesis.
      \item \textbf{Adaptive Engine}:  By coupling the learner model with the
            semantic knowledge base, the adaptive engine generates customized
            learning pathways, selects appropriate practice items, and delivers
            real‑time feedback\footnote{Green yellow red light}.  It also recommends prerequisite content when a
            knowledge gap is detected, thus scaffolding the student’s learning
            experience.
\end{enumerate}
Through this triadic architecture, ALeA offers a coherent environment in which
semantic knowledge representation, dynamic learner modelling, and adaptive
instruction converge to enhance both teaching efficacy and student outcomes.

